% ============================================================================
\chapter{Differentiable Gate Selection}
\label{ch:diff-gates}
% ============================================================================

This chapter introduces the Logic Gate Network (LGN)~\cite{petersen2022difflogic}, a neural network
where each node is a 2-input logic gate. The key innovation is a
training method that uses gradient descent to \emph{select} which of
the 16 possible gate types each node should implement. After training,
the continuous selection mechanism is removed, leaving a pure Boolean
circuit.


% ----------------------------------------------------------------------------
\section{The 16 Two-Input Boolean Functions}
\label{sec:sixteen-gates}
% ----------------------------------------------------------------------------

As we saw in \cref{sec:truth-tables}, two binary inputs $(a, b)$
have $2^2 = 4$ input combinations, and each combination maps to 0
or~1. This gives $2^4 = 16$ possible functions. We list all 16:

\begin{table}[htbp]
\centering
\caption{All 16 two-input Boolean functions, ordered by their 4-bit
  truth table code. The code lists outputs for inputs $(0,0),
  (0,1), (1,0), (1,1)$.}
\label{tab:all-16}
\begin{tabular}{@{}clcccc@{}}
  \toprule
  \textbf{Index} & \textbf{Name} & $(0,0)$ & $(0,1)$ & $(1,0)$ & $(1,1)$ \\
  \midrule
  0  & FALSE (constant 0) & 0 & 0 & 0 & 0 \\
  1  & AND              & 0 & 0 & 0 & 1 \\
  2  & $A \wedge \bar{B}$ (inhibit B) & 0 & 0 & 1 & 0 \\
  3  & $A$ (project A)    & 0 & 0 & 1 & 1 \\
  4  & $\bar{A} \wedge B$ (inhibit A) & 0 & 1 & 0 & 0 \\
  5  & $B$ (project B)    & 0 & 1 & 0 & 1 \\
  6  & XOR              & 0 & 1 & 1 & 0 \\
  7  & OR               & 0 & 1 & 1 & 1 \\
  8  & NOR              & 1 & 0 & 0 & 0 \\
  9  & XNOR             & 1 & 0 & 0 & 1 \\
  10 & $\bar{B}$ (NOT B)  & 1 & 0 & 1 & 0 \\
  11 & $A \vee \bar{B}$ (imply A) & 1 & 0 & 1 & 1 \\
  12 & $\bar{A}$ (NOT A)  & 1 & 1 & 0 & 0 \\
  13 & $\bar{A} \vee B$ (imply B) & 1 & 1 & 0 & 1 \\
  14 & NAND             & 1 & 1 & 1 & 0 \\
  15 & TRUE (constant 1)  & 1 & 1 & 1 & 1 \\
  \bottomrule
\end{tabular}
\end{table}

Each function is uniquely identified by its 4-bit truth table code.
We denote the truth table of gate type $g$ as $\TruthTable(g) \in
\{0,1\}^4$, indexed by the input combination.


% ----------------------------------------------------------------------------
\section{The Gate Selection Problem}
\label{sec:gate-selection}
% ----------------------------------------------------------------------------

The Logic Gate Network~\cite{petersen2022difflogic,petersen2024difflogic} consists of layers of nodes. Each node takes
two binary inputs and produces one binary output. The question is:
which of the 16 gate types should each node implement?

If we fix the gate types, evaluation is trivial---just apply the logic
operations. The challenge is \emph{learning} the gate types from data.

\begin{keyinsight}
A logic gate is a discrete object: it is AND or it is OR; there is no
``halfway between AND and OR.'' Gradient descent cannot directly
optimise a discrete choice. The LGN solves this by replacing the
discrete choice with a continuous probability distribution over gate
types, optimising this distribution with gradients, and then
discretising after training.
\end{keyinsight}


% ----------------------------------------------------------------------------
\section{The Continuous Relaxation}
\label{sec:relaxation}
% ----------------------------------------------------------------------------

For each node $i$ in the network, we maintain a vector of 16
\concept{logits}:
\begin{equation}
  \bvec{z}_i = (z_i^{(0)}, z_i^{(1)}, \ldots, z_i^{(15)})
    \in \mathbb{R}^{16}
  \label{eq:logits}
\end{equation}

These logits are converted to a probability distribution using the
\concept{softmax} function:
\begin{equation}
  p_i^{(g)} = \frac{\exp(z_i^{(g)} / \tau)}
    {\sum_{g'=0}^{15} \exp(z_i^{(g')} / \tau)}
  \label{eq:softmax-gate}
\end{equation}

where $\tau > 0$ is the \concept{temperature} parameter.

The node's output is then a \emph{weighted average} of all 16 gate
outputs:
\begin{equation}
  \hat{y}_i = \sum_{g=0}^{15} p_i^{(g)} \cdot \TruthTable(g)[a_i, b_i]
  \label{eq:relaxed-output}
\end{equation}

where $a_i, b_i \in \{0, 1\}$ are the node's two inputs, and
$\TruthTable(g)[a_i, b_i]$ is the output of gate type $g$ on those
inputs.

\begin{example}
Suppose a node has inputs $a = 1, b = 0$. The truth table outputs for
this input combination are: AND gives 0, OR gives 1, XOR gives 1,
NAND gives 1, etc. If the probabilities are
$p^{(\text{AND})} = 0.6$ and $p^{(\text{OR})} = 0.3$ and
$p^{(\text{others})} = 0.1$, then:
\[
  \hat{y} = 0.6 \times 0 + 0.3 \times 1 + 0.1 \times (\text{others})
  = 0.3 + \cdots \approx 0.4
\]
The output is continuous (between 0 and 1), not binary. This is the
price of differentiability.
\end{example}


% ----------------------------------------------------------------------------
\section{Temperature Annealing}
\label{sec:annealing}
% ----------------------------------------------------------------------------

The temperature $\tau$ controls how ``soft'' the gate selection is:

\begin{itemize}
  \item $\tau \to \infty$: All gate types have equal probability
    ($p^{(g)} = 1/16$). The output is a uniform average. Maximum
    exploration, minimum commitment.
  \item $\tau \to 0$: The probability concentrates on a single gate
    type (the one with the largest logit). The output approaches a
    hard Boolean value. Maximum commitment, zero exploration.
\end{itemize}

During training, we anneal $\tau$ from a high value to a low value:

\begin{equation}
  \tau(t) = \tau_{\text{start}} \cdot
    \left(\frac{\tau_{\text{end}}}{\tau_{\text{start}}}\right)^{t / t_{\max}}
  \label{eq:anneal}
\end{equation}

where $t$ is the training step and $t_{\max}$ is the total number of
steps.

\begin{figure}[htbp]
\centering
\begin{tikzpicture}[>=Stealth]
  \begin{axis}[
    width=10cm, height=5cm,
    xlabel={Training progress (\%)}, ylabel={Temperature $\tau$},
    xmin=0, xmax=100, ymin=0, ymax=5,
    grid=major, grid style={bitgray!20},
    every axis label/.style={font=\small\sffamily},
    tick label style={font=\small}
  ]
    \addplot[thick, bitblue, smooth, domain=0:100, samples=50]
      {5 * (0.01/5)^(x/100)};

    \node[font=\scriptsize\sffamily, bitblue] at (axis cs:15, 4)
      {explore};
    \node[font=\scriptsize\sffamily, bitblue] at (axis cs:85, 0.8)
      {commit};
  \end{axis}
\end{tikzpicture}
\caption{Temperature annealing schedule. The temperature starts high
  (soft selection, exploration) and decays to near zero (hard
  selection, commitment). The exponential decay ensures a gradual
  transition.}
\label{fig:annealing}
\end{figure}

\begin{analogy}
Temperature annealing is like a hiring decision. Early in the process,
you keep many candidates under consideration (high temperature). As
interviews proceed, you narrow the field. Eventually, you commit to
one candidate (low temperature). The gradual narrowing avoids premature
decisions based on limited information.
\end{analogy}


% ----------------------------------------------------------------------------
\section{Why This Works: The Gradient Flow}
\label{sec:gradient-flow}
% ----------------------------------------------------------------------------

The relaxed output (\cref{eq:relaxed-output}) is differentiable with
respect to the logits $\bvec{z}_i$. The gradient of the loss with
respect to logit $z_i^{(g)}$ tells us: ``if we increase the
probability of gate type $g$ at node $i$, does the loss go up or
down?''

Specifically, if $\mathcal{L}$ is the loss:
\begin{equation}
  \frac{\partial \mathcal{L}}{\partial z_i^{(g)}} =
    \frac{\partial \mathcal{L}}{\partial \hat{y}_i} \cdot
    \frac{\partial \hat{y}_i}{\partial z_i^{(g)}}
  \label{eq:gate-grad}
\end{equation}

The second factor involves the softmax derivative and the truth table
values. Standard backpropagation through the network computes the
first factor.

\begin{engineernote}
The implementation is straightforward in any automatic differentiation
framework. For each node, the forward pass computes:
\begin{enumerate}
  \item Look up the 4 truth table outputs for all 16 gate types (a
    constant $16 \times 4$ matrix).
  \item Index into this matrix using the binary inputs $a_i, b_i$ to
    get a 16-element vector of gate outputs.
  \item Compute softmax of logits to get probabilities.
  \item Dot product of probabilities and gate outputs to get
    $\hat{y}_i$.
\end{enumerate}
The backward pass is handled automatically by the framework.
\end{engineernote}


% ----------------------------------------------------------------------------
\section{Discretisation: From Soft to Hard}
\label{sec:discretisation}
% ----------------------------------------------------------------------------

After training with fully annealed temperature, each node's
probability distribution is sharply peaked at one gate type. We
``harden'' the network by selecting the most probable gate:

\begin{equation}
  g_i^* = \arg\max_g \; z_i^{(g)}
  \label{eq:hardening}
\end{equation}

The logits and softmax machinery are then discarded. What remains is
a fixed Boolean circuit: each node is a specific logic gate with two
inputs and one output.

\begin{warning}
Discretisation introduces a small accuracy drop because the
continuous weighted average is replaced by a single gate. The drop
is typically 0.1--0.5\% on MNIST~\cite{petersen2024difflogic}. If the annealing schedule is too
aggressive (temperature drops too fast), the drop can be larger
because nodes commit before finding good gate types.
\end{warning}


% ----------------------------------------------------------------------------
\section{Network Topology}
\label{sec:lgn-topology}
% ----------------------------------------------------------------------------

The LGN arranges nodes in layers. Each node in layer $l$ takes two
inputs from layer $l-1$ (or from the input). The connectivity can be:

\begin{description}[leftmargin=1em]
  \item[Random sparse:] Each node's two inputs are randomly chosen
    from the previous layer. This is the approach
    in~\cite{petersen2022difflogic}. Despite the randomness, the
    learned gate types compensate, and the network achieves
    surprisingly good accuracy.

  \item[Fully connected (impractical):] Every pair of nodes in the
    previous layer could be an input. This has quadratic cost and is
    not used.

  \item[Structured:] Inputs are chosen based on some fixed pattern
    (e.g., adjacent clause outputs). This can encode domain knowledge
    but limits flexibility.
\end{description}

For our hybrid architecture, we use random sparse connectivity.
Each layer has a fixed number of nodes, and each node's two inputs
are independently sampled (with replacement) from the previous layer's
outputs.

\begin{engineernote}
A typical LGN for MNIST~\cite{petersen2024difflogic}: 4--6 layers, each with 2000--5000 nodes,
random connectivity. The final layer's outputs are aggregated
(e.g., by popcount or a small voting circuit) to produce the
classification decision. Total gate count: 10,000--30,000.
At inference, each gate is a single table lookup (4 entries) or a
bitwise operation.
\end{engineernote}


% ----------------------------------------------------------------------------
\section{The Straight-Through Estimator Alternative}
\label{sec:ste}
% ----------------------------------------------------------------------------

An alternative to the softmax relaxation is the \concept{straight-through
estimator} (STE)~\cite{bengio2013ste}. During the forward pass, the node
commits to a hard gate type (the argmax). During the backward pass, the
gradient is computed \emph{as if} the soft relaxation were used.

This avoids the accuracy gap between soft and hard evaluation, but
introduces a bias in the gradient (since the backward pass does not
match the forward pass). In practice, both approaches work well for
LGNs; the softmax relaxation is more principled, while STE can be
faster.


% ----------------------------------------------------------------------------
\section{Summary}
\label{sec:ch8-summary}
% ----------------------------------------------------------------------------

The Logic Gate Network makes Boolean circuits learnable through three
ideas:
\begin{enumerate}
  \item \textbf{Enumerate all gates:} There are only 16 two-input
    Boolean functions. This finite set enables exhaustive
    consideration.
  \item \textbf{Softmax relaxation:} A probability distribution over
    gate types makes the output continuous and differentiable.
    Gradient descent optimises the logits.
  \item \textbf{Temperature annealing:} Gradually reducing the
    temperature transitions from exploration (all gates equally
    likely) to commitment (one gate selected).
\end{enumerate}

After training and discretisation, the LGN is a pure Boolean circuit
with zero continuous-valued operations. The next chapter describes
how to train this network on Tsetlin Machine clause outputs.
